\documentclass{article}
\usepackage{amsmath, amssymb, amsthm}
\usepackage{hyperref}
\usepackage{geometry}
\geometry{margin=1in}

\title{Erd\H{o}s Problem \#1160}
\date{Last updated: 2026-01-23}
\author{Source: erdosproblems.com}

\begin{document}
\maketitle

\noindent\textbf{Status:} OPEN \quad \textbf{No prize}

\noindent\textbf{Tags:} group theory

\medskip
\noindent\textbf{Problem Statement:}

Let $g(n)$ denote the number of groups of order $n$. If $n\leq 2^m$ then $g(n)\leq g(2^m)$.

\bigskip
\noindent\textbf{Remarks:}

This is listed as an open problem (Question 22.16) in \cite{BNV07}, which reports it as a 'quite natural conjecture, whose origin we have been unable to trace satisfactorily. We have heard it attributed at various times to various people, such as Paul Erd\H{o}s and Graham Higman.'

Question 22.18 of \cite{BNV07} suggests the even stronger conjecture\[\sum_{n<2^m}g(n) \leq g(2^m)\]for all sufficiently large $m$ (perhaps even as soon as $m\geq 7$).

Pantelidakis \cite{Pa03} proved that the original conjecture is true if $n$ is odd and $m\geq 3619$.

\bigskip
\noindent\textbf{References:}

[BNV07] Blackburn, Simon R. and Neumann, Peter M. and Venkataraman,
Geetha, Enumeration of finite groups. (2007), xii+281.
 
 [Pa03] I. Pantelidakis, On the Number of Non-isomorphic Groups of the Same Order. DPhil Thesis, University of Oxford (2003).

\bigskip
\noindent\small{Source: \url{https://www.erdosproblems.com/1160}}
\end{document}
